\documentclass[12pt,a4paper]{article}
\usepackage{fontspec}
\usepackage[italian]{babel}
\usepackage{tikz}
\usepackage{pgfplots}
\usepackage{xstring}
\usepackage{pgfplotstable}
\usepgfplotslibrary{polar}
\pgfplotsset{compat=1.18}

% Librerie TikZ necessarie
\usetikzlibrary{shapes,arrows,calc,patterns,decorations.pathreplacing}
\usetikzlibrary{shadows,backgrounds,fit,trees,mindmap}
\usetikzlibrary{charts,graphs,er,folding,calendar,matrix}
\usetikzlibrary{intersections,through,datavisualization}
\usetikzlibrary{decorations.text,calligraphy}
\usetikzlibrary{positioning}

% Definizione colori GIST
\definecolor{gistblue}{RGB}{52,152,219}      % Physical
\definecolor{gistgreen}{RGB}{46,204,113}     % Architectural  
\definecolor{gistred}{RGB}{231,76,60}        % Security
\definecolor{gistyellow}{RGB}{241,196,15}    % Compliance
\definecolor{gistgray}{RGB}{149,165,166}     % Neutral

% Stili comuni
\tikzset{
    component/.style={rectangle,rounded corners,minimum width=3cm,minimum height=1.5cm,font=\sffamily},
    physical/.style={component,fill=gistblue!80,text=white},
    architectural/.style={component,fill=gistgreen!80,text=white},
    security/.style={component,fill=gistred!80,text=white},
    compliance/.style={component,fill=gistyellow!80,text=black},
    arrow/.style={->,>=stealth,thick},
    dashedarrow/.style={->,>=stealth,thick,dashed}
}

% Map 'status' token to a color
\newcommand{\statuscolor}[1]{%
  \IfStrEq{##1}{done}{gistgreen}{%
  \IfStrEq{##1}{current}{gistyellow}{gray}}}
\begin{document}

\title{Nuove Figure per Nuova Tesi}
\begin{figure}[htbp]
\centering
\begin{tikzpicture}
    \begin{polaraxis}[
        title={\textbf{Gap tra Ricerca e Implementazione nella GDO}},
        title style={font=\Large},
        width=12cm,
        height=12cm,
        ylabel near ticks,
        ylabel style={font=\small},
        xtick={0,60,120,180,240,300},
        xticklabels={
            Approccio\\Olistico,
            Modelli\\Economici,
            Vincoli\\Operativi,
            Specificità\\Settoriale,
            Integrazione\\Normativa,
            Applicabilità\\Pratica
        },
        xticklabel style={align=center,font=\small},
        ytick={0,20,40,60,80,100},
        ymin=0, ymax=100,
        ylabel={Copertura (\%)},
        grid=both,
        major grid style={gray!30},
        minor grid style={gray!10},
        legend pos=south east,
        legend style={font=\footnotesize}
    ]
    
    % Area Ricerca Accademica
    \addplot[thick,gistblue,fill=gistblue!30,opacity=0.7] 
        coordinates {
            (0,45) (60,35) (120,38) (180,42) (240,40) (300,48) (360,45)
        };
    \addlegendentry{Ricerca Accademica}
    
    % Area Requisiti Reali
    \addplot[thick,gistgreen,fill=gistgreen!30,opacity=0.7] 
        coordinates {
            (0,87) (60,92) (120,89) (180,95) (240,88) (300,90) (360,87)
        };
    \addlegendentry{Requisiti GDO}
    
    % Area Gap (pattern)
    \addplot[pattern=north east lines,pattern color=gistred,opacity=0.5] 
        coordinates {
            (0,87) (60,92) (120,89) (180,95) (240,88) (300,90) (360,87)
        } --
        coordinates {
            (360,45) (300,48) (240,40) (180,42) (120,38) (60,35) (0,45)
        };
    \addlegendentry{Gap Identificato}
    
    % Annotazioni
    \node[font=\footnotesize,fill=white,inner sep=2pt] at (axis cs:60,65) {Gap: 62\%};
    \node[font=\footnotesize,fill=white,inner sep=2pt] at (axis cs:300,69) {Gap: 38\%};
    
    \end{polaraxis}
    
    % Box con media gap
    \node[draw,fill=gistgray!20,anchor=north west,font=\small] at (7,-0.5) {
        \begin{tabular}{l}
        \textbf{Gap Medio:} 51.7\%\\
        \textbf{Max Gap:} Modelli Economici\\
        \textbf{Min Gap:} Approccio Olistico
        \end{tabular}
    };
\end{tikzpicture}
\caption{Gap identificato tra ricerca accademica e requisiti reali del settore GDO}
\label{fig:gap-ricerca}
\end{figure}
\begin{figure}[htbp]
\begin{tikzpicture}[node distance=2cm]
    % Livello 1: Componenti Base
    \node[physical] (phys) at (0,0) {\parbox[c][1.5cm][c]{3cm}{\centering\textbf{Physical}\\18\%}};
    \node[architectural,right=of phys] (arch) {\parbox[c][1.5cm][c]{3cm}{\centering\textbf{Architectural}\\32\%}};
    \node[security,right=of arch] (sec) {\parbox[c][1.5cm][c]{3cm}{\centering\textbf{Security}\\28\%}};
    \node[compliance,right=of sec] (comp) {\parbox[c][1.5cm][c]{3cm}{\centering\textbf{Compliance}\\22\%}};
    
    % Livello 2: Elaborazione
    \node[draw,diamond,aspect=2,above=3cm of arch,fill=gray!20] (decision) {Modalità};
    
    % Percorso Balanced (sinistra)
    \node[draw,rectangle,left=2cm of decision,fill=blue!20,align=center] (balanced) 
        {\(\sum_{i}(w_i \times C_i)\)\\[0.3cm]\textit{Balanced}};
    
    % Percorso Critical (destra)
    \node[draw,rectangle,right=2cm of decision,fill=red!20,align=center] (critical) 
        {\(\prod_{i}(C_i^{w_i})\)\\[0.3cm]\textit{Critical}};
    
    % Livello 3: Output
    \node[draw,rectangle,above=5cm of decision,minimum width=4cm,minimum height=2cm,fill=green!20,align=center] (output) 
        {\textbf{GIST Score}\\0--100};
    
    % Moltiplicatori laterali
    \node[draw,circle,left=of output,fill=yellow!20,align=center] (kgdo) {\(K_{GDO}\)\\1.0--1.5};
    \node[draw,circle,right=of output,fill=orange!20,align=center] (innov) {\((1+I)\)\\1.0--1.35};
    
    % Frecce dal livello 1 al decision point
    \draw[arrow] (phys.north) -- ++(0,1.5) -| (decision);
    \draw[arrow] (arch.north) -- (decision);
    \draw[arrow] (sec.north) -- ++(0,1.5) -| (decision);
    \draw[arrow] (comp.north) -- ++(0,1.5) -| (decision);
    
    % Frecce da decision a modalità
    \draw[arrow] (decision) -- (balanced);
    \draw[arrow] (decision) -- (critical);
    
    % Frecce da modalità a output
    \draw[arrow] (balanced.north) -- ++(0,1) -| (output);
    \draw[arrow] (critical.north) -- ++(0,1) -| (output);
    
    % Frecce moltiplicatori
    \draw[arrow] (kgdo) -- (output);
    \draw[arrow] (innov) -- (output);
    
    % Etichette
    \node[font=\footnotesize,below=0.1cm of phys,align=center] {Infrastruttura\\Fisica};
    \node[font=\footnotesize,below=0.1cm of arch,align=center] {Maturità\\Architetturale};
    \node[font=\footnotesize,below=0.1cm of sec,align=center] {Postura\\Sicurezza};
    \node[font=\footnotesize,below=0.1cm of comp,align=center] {Integrazione\\Compliance};
    
    % Switch indicatore
    \draw[thick,dotted] ($(decision.west)-(0.5,0)$) -- ($(decision.east)+(0.5,0)$);
    \node[font=\tiny,above=0.1cm of decision] {Switch};
\end{tikzpicture}
\end{tikzpicture}
\caption{Architettura del Framework GIST con doppia formulazione}
\label{fig:gist-architecture}
\end{figure}
\begin{figure}[htbp]
\centering
\begin{tikzpicture}
    % Definizione del cerchio principale
    \def\radius{4cm}
    \coordinate (center) at (0,0);
    
    % Componenti principali (in cerchio)
    \node[draw,circle,fill=gistblue!80,text=white,minimum size=2cm] (threat) at (90:\radius) {
        \textbf{Threat}\\
        \textbf{Intelligence}
    };
    
    \node[draw,circle,fill=gistgreen!80,text=white,minimum size=2cm] (risk) at (0:\radius) {
        \textbf{Risk}\\
        \textbf{Assessment}
    };
    
    \node[draw,circle,fill=gistred!80,text=white,minimum size=2cm] (control) at (270:\radius) {
        \textbf{Control}\\
        \textbf{Selection}
    };
    
    \node[draw,circle,fill=gistyellow!80,text=black,minimum size=2cm] (impl) at (180:\radius) {
        \textbf{Implementation}\\
    };
    
    % Inner loop (tattico)
    \draw[arrow,thick,gistblue] (threat) to[bend left=20] node[midway,right,font=\small] {Prioritize} (risk);
    \draw[arrow,thick,gistgreen] (risk) to[bend left=20] node[midway,below,font=\small] {Optimize} (control);
    \draw[arrow,thick,gistred] (control) to[bend left=20] node[midway,left,font=\small] {Deploy} (impl);
    \draw[arrow,thick,gistyellow] (impl) to[bend left=20] node[midway,above,font=\small] {Monitor} (threat);
    
    % Outer loop (strategico)
    \draw[dashedarrow,thick,gray] (threat) to[bend right=45] node[midway,left,font=\small] {Strategic\\Review} (control);
    \draw[dashedarrow,thick,gray] (risk) to[bend right=45] node[midway,right,font=\small] {Policy\\Update} (impl);
    
    % Centro con metriche
    \node[draw,circle,fill=white,minimum size=3cm] at (center) {
        \begin{tabular}{c}
        \textbf{Real-Time Metrics}\\[0.2cm]
        ASSA: 1,631\\
        Incidents: 4.2/month\\
        MTTD: 24h\\
        Compliance: 94\%
        \end{tabular}
    };
    
    % Timeline indicators
    \node[font=\footnotesize,gray] at ($(threat)+(0,1.2)$) {Continuous};
    \node[font=\footnotesize,gray] at ($(risk)+(1.2,0)$) {Weekly};
    \node[font=\footnotesize,gray] at ($(control)+(0,-1.2)$) {Monthly};
    \node[font=\footnotesize,gray] at ($(impl)+(-1.2,0)$) {Quarterly};
    
    % Legenda
    \node[anchor=north west,font=\footnotesize] at (-5,-5) {
        \textbf{Cicli di Feedback:}\\
        \textcolor{gray}{---} Loop Strategico (mensile)\\
        \textcolor{black}{$\rightarrow$} Loop Tattico (24h)
    };
\end{tikzpicture}
\caption{Framework integrato di sicurezza GDO con cicli di feedback}
\label{fig:security-framework}
\end{figure}

\begin{figure}[htbp]
\centering
\begin{tikzpicture}
\begin{semilogxaxis}[
    title={\textbf{Correlazione tra Investimenti e Availability}},
    xlabel={Investimento Infrastrutturale (€/negozio)},
    ylabel={Availability (\%)},
    xmin=1000, xmax=50000,
    ymin=99.0, ymax=100.0,
    width=14cm,
    height=10cm,
    grid=both,
    minor grid style={gray!10},
    major grid style={gray!30},
    legend pos=south east,
    legend style={font=\footnotesize}
]

% Nuvola di punti Monte Carlo
\addplot[
    only marks,
    mark=*,
    mark size=0.5pt,
    gray!50,
    opacity=0.3
] table[x=investment,y=availability] {
    % Qui andrebbero i dati reali
    % Per esempio dimostrativo:
    investment availability
    1500 99.12
    2000 99.23
    2500 99.31
    3000 99.38
    4000 99.48
    5000 99.56
    7000 99.68
    10000 99.78
    15000 99.86
    20000 99.91
    30000 99.94
    40000 99.96
};

% Punti reali
\addplot[
    only marks,
    mark=*,
    mark size=2pt,
    gistblue
] coordinates {
    (8700,99.82) (12300,99.96) (11200,99.94)
};

% Curve di fitting
\addplot[
    domain=1000:50000,
    samples=100,
    thick,
    black
] {99.0 + 0.7*ln(x/1000)};
\addlegendentry{Regressione ($R^2=0.84$)}

% Bande di confidenza
\addplot[
    domain=1000:50000,
    samples=100,
    thick,
    gistgreen,
    dashed,
    name path=upper
] {99.0 + 0.7*ln(x/1000) + 0.1};

\addplot[
    domain=1000:50000,
    samples=100,
    thick,
    gistgreen,
    dashed,
    name path=lower
] {99.0 + 0.7*ln(x/1000) - 0.1};

\addplot[gistgreen!20,opacity=0.3] fill between[of=upper and lower];
\addlegendentry{IC 95\%}

% Target line
\addplot[
    domain=1000:50000,
    samples=2,
    thick,
    gistred,
    dashed
] {99.95};
\addlegendentry{Target 99.95\%}

% Zone annotate
\node[draw,fill=red!20,opacity=0.7,anchor=south west] at (axis cs:1000,99.8) {
    \begin{tabular}{c}
    \textbf{Underinvestment}\\
    High Risk
    \end{tabular}
};

\node[draw,fill=green!20,opacity=0.7] at (axis cs:10000,99.5) {
    \begin{tabular}{c}
    \textbf{Optimal Zone}\\
    Best ROI
    \end{tabular}
};

\node[draw,fill=yellow!20,opacity=0.7,anchor=south east] at (axis cs:40000,99.3) {
    \begin{tabular}{c}
    \textbf{Diminishing}\\
    \textbf{Returns}
    \end{tabular}
};

% Annotazioni punti chiave
\node[pin={[pin edge={gistblue,thick}]90:Break-even}] at (axis cs:8700,99.82) {};
\node[pin={[pin edge={gistblue,thick}]45:Sweet spot}] at (axis cs:12300,99.96) {};

\end{semilogxaxis}
\end{tikzpicture}
\caption{Relazione tra investimenti infrastrutturali e disponibilità del servizio}
\label{fig:investment-availability}
\end{figure}
\begin{figure}[htbp]
\centering
\begin{tikzpicture}[scale=0.9]
    % Centro - GIST Score
    \node[draw,circle,minimum size=3cm,fill=white] (center) at (0,0) {
        \begin{tabular}{c}
        \textbf{GIST Score}\\[0.3cm]
        \huge\textbf{76.8}\\[0.2cm]
        \textit{Maturo}
        \end{tabular}
    };
    
    % Gauge meter background
    \draw[thick,gray!30] (0,0) circle (1.8cm);
    \draw[thick,gistgreen,line width=4pt] ([shift=(140:1.8cm)]0,0) arc (140:40:1.8cm);
    
    % Quadranti
    \def\quadradius{5}
    \def\compradius{2}
    
    % Physical (Nord)
    \node[draw,circle,fill=gistblue!80,text=white,minimum size=\compradius cm] (physical) at (90:\quadradius) {
        \begin{tabular}{c}
        \textbf{Physical}\\
        68/100
        \end{tabular}
    };
    
    % Architectural (Est)
    \node[draw,circle,fill=gistgreen!80,text=white,minimum size=\compradius cm] (arch) at (0:\quadradius) {
        \begin{tabular}{c}
        \textbf{Architectural}\\
        82/100
        \end{tabular}
    };
    
    % Security (Sud)
    \node[draw,circle,fill=gistred!80,text=white,minimum size=\compradius cm] (security) at (270:\quadradius) {
        \begin{tabular}{c}
        \textbf{Security}\\
        74/100
        \end{tabular}
    };
    
    % Compliance (Ovest)
    \node[draw,circle,fill=gistyellow!80,text=black,minimum size=\compradius cm] (compliance) at (180:\quadradius) {
        \begin{tabular}{c}
        \textbf{Compliance}\\
        79/100
        \end{tabular}
    };
    
    % Mini-indicatori per componente
    \foreach \comp/\angle/\score in {physical/90/68,arch/0/82,security/270/74,compliance/180/79} {
        \draw[thick,gray!30] ([shift=(\angle:\quadradius)]0,0) ++(0.2,0.8) -- ++(1,0);
        \draw[thick,gistgreen,line width=2pt] ([shift=(\angle:\quadradius)]0,0) ++(0.2,0.8) -- ++({\score/100}cm,0);
    }
    
    % Connessioni sinergiche
    \draw[thick,gistgreen!50,<->,>=stealth] (physical) to[bend left=30] node[midway,font=\footnotesize] {+27\%} (arch);
    \draw[thick,gistgreen!50,<->,>=stealth] (arch) to[bend left=30] node[midway,font=\footnotesize] {+34\%} (security);
    \draw[thick,gistgreen!50,<->,>=stealth] (security) to[bend left=30] node[midway,font=\footnotesize] {+41\%} (compliance);
    \draw[thick,gistgreen!50,<->,>=stealth] (compliance) to[bend left=30] node[midway,font=\footnotesize] {+23\%} (physical);
    
    % Anello esterno - Timeline
    \draw[thick,gray!20] (0,0) circle (7cm);
    
    % Milestone
    \foreach \angle/\label/\status in {
        45/{M1: Kick-off}/done,
        0/{M3: Cloud Wave 1}/done,
        315/{M6: SD-WAN}/done,
        270/{M9: Zero Trust}/current,
        225/{M12: Multi-cloud}/planned,
        180/{M18: Automation}/planned,
        135/{M24: Steady State}/planned
    }%
\pgfplotsforeachungrouped \angle/\label/\status in {
        \node[circle,fill=\statuscolor{\status},minimum size=0.3cm] at (\angle:7) {};
        \node[font=\tiny,rotate=\angle,anchor=south,inner sep=1pt] at (\angle:7.5) {\label};
    }
    
    % Indicatore progresso
    \draw[thick,gistgreen,line width=3pt,->,>=stealth] (45:7) arc (45:270:7);
    
    % Legenda sinergie
    \node[anchor=north,font=\footnotesize] at (0,-8) {
        \textbf{Effetto Sistema:} +52\% oltre somma lineare componenti
    };
\end{tikzpicture}
\caption{Framework GIST integrato - Vista sintetica con sinergie}
\label{fig:gist-integrated}
\end{figure}
% Waterfall Chart per Figura 2.2
\begin{figure}[htbp]
\centering
\begin{tikzpicture}
    \begin{axis}[
        ybar,
        width=14cm,
        height=10cm,
        ymin=0,
        ymax=3000,
        ylabel={ASSA Score},
        symbolic x coords={Baseline,Microseg,Edge Isol,Traffic Insp,Monitoring,Altri,Finale},
        xtick=data,
        x tick label style={rotate=45,anchor=east,font=\small},
        nodes near coords,
        every node near coord/.append style={font=\footnotesize},
        grid=major,
        grid style={gray!30}
    ]
    
    % Valori cumulativi per waterfall
    \addplot[fill=gistred!80] coordinates {
        (Baseline,2847)
        (Microseg,0)
        (Edge Isol,0)
        (Traffic Insp,0)
        (Monitoring,0)
        (Altri,0)
        (Finale,1631)
    };
    
    % Riduzioni (negative bars)
    \addplot[fill=gistblue!80] coordinates {
        (Baseline,0)
        (Microseg,-888)
        (Edge Isol,-687)
        (Traffic Insp,-524)
        (Monitoring,-371)
        (Altri,-109)
        (Finale,0)
    };
    
    % Linee di connessione
    \draw[dashed,thick] (axis cs:Baseline,2847) -- (axis cs:Microseg,1959);
    \draw[dashed,thick] (axis cs:Microseg,1959) -- (axis cs:{Edge Isol},1272);
    % ... continua per altre connessioni
    
    \end{axis}
\end{tikzpicture}
\caption{Decomposizione della riduzione ASSA per componente}
\label{fig:assa-waterfall}
\end{figure}
\end{document}
% End of document